\section{Introduction}


The ATLAS physics programme relies on an efficient trigger system~\cite{TRIG-2016-01}, allowing the experiment to record highly signal-enriched subsets of collision events produced by the Large Hadron Collider (LHC) at CERN~\cite{Evans_2007}. Electron triggers~\cite{TRIG-2018-05} are important for searches for new particles\cite{EXOT-2019-13, EXOT-2018-55}, measurements of Standard Model cross-sections and properties of fundamental particles such as the Higgs~\cite{HIGG-2012-27,HIGG-2016-33}, $Z$ and $W$ bosons.

This paper describes the performance improvements achieved by the ATLAS trigger system thanks to the introduction of a neural network based algorithm, the Neural-Ringer, to identify electrons, while rejecting QCD backgrounds efficiently. The Neural-Ringer became the baseline algorithm for selecting isolated electrons online by the ATLAS trigger system in the 2017 proton--proton collision at $\sqs = 13$~TeV in LHC \RunTwo.  


The ATLAS detector and the trigger system are described in section~\ref{sec:atlas}. Electron identification, both offline and online, is described in section \ref{sec:electronID}. 
Full details of the \rnn electron identification, including the training strategies, are given in section~\ref{sec:neuralringer}. Online performance of the \rnn and its comparison with respect to the previous trigger baseline algorithm are described in section~\ref{sec:operation}. A detailed statistical analysis of the \rnn electron results for electron identification is presented in section~\ref{sec:off_ana}. Finally, the conclusions and prospects are presented in section~\ref{sec:conclusion}.
% 
% 
%  Electron selection starts with a fast calorimeter-reconstruction step, which has two implementations: the original cut-based selection and the new neural-network (NN) approach using the Neural-Ringer algorithm presented in this paper. This is followed by a fast track-reconstruction step, which also places requirements on variables that indicate the quality of proximity matching between the electron candidate positions reconstructed from the track and the calorimetric measurements. The HLT tries to select particle candidates at each step; if it fails at a certain step, subsequent steps are not executed. This is essential in order to reduce the CPU resources and time needed by the HLT to reconstruct the event and make a trigger decision. 





